\cotefile{40}

\scanpage{2}

% archive sheet number: 1

On se place dans la catégorie duale de la catégorie des algèbres $B$ (commutatives avec unité) sur l'anneau comm. $A$.

% slanted
\margin{\uline{Calculs \inserted{affines} sur} des fibrés \uline{principaux} \uline{affines}}

Une telle $B$ est un groupe de la catégorie si \unsure{c'est|l'on} une \uline{$A$-algèbre commutative unitaire} s'est donné un \uline{hom. d'algèbres}

\[
\unsure{p|\rho}_{B}\colon B\longrightarrow B\otimes_A B
\]

(``application diagonale'') qui soit ``\uline{associative}'', pour laquelle il \uline{existe} une \uline{augmentation} \inserted{\uline{d'algèbre}} $\varepsilon:B\to A$ \deleted{telle que} telle que

\[
\delta_B x=p_Bx-\deleted{\unsure{( |x}}(x\otimes1+1\otimes x)\in\unsure{J|I}\otimes J
\]

(où $J=\operatorname{Ker}\varepsilon$) --- cette augmentation est alors unique --- et un \uline{hom d'algèbres} $x\longrightarrow\check{x}$ satisfaisant à la condition

\[
x\in B,\qquad p_Bx=\sum y_i z_i\quad\Longrightarrow\quad \sum\check{y}_i z_i=\sum y_i\check{z}_i=\varepsilon(x).1
\]

--- cet hom d'algèbres est alors unique, et est une involution. De plus on aura

\[
\varepsilon(\check{x})=\varepsilon(x)
\]

\[
\text{\deleted{\unsure{et}} si}\qquad p_B\unsure{'}x=\sum y_i\otimes z_i,\quad\text{alors}\quad p_B\check{x}=\sum_i\check{z}_i\otimes\check{y}_i
\]

\uline{Si} $B$ est libre de rang fini sur $B$, alors\deleted{\unsure{s}} ces structures se reflètent directement dans $U$ par la donnée \deleted{d'une \uline{loi d'algèbre} (\unsure{non|i.e.} nécessairement commutative) \uline{associative unitaire} sur $A$, \inserted{dont la duale est une} telle qu'il existe un \unsure{hom} \unsure{$x\to\check{x}$|$x\mapsto\check{x}$}}
